\section{模型校验、评价与推广}\label{sec:eval}

\subsection{模型与程序的自校验}

所有结论都由确定性程序给出，程序在每次运行时自动执行一组互相独立的校验，结果写入结论文件。
表~\ref{tab:checks} 汇总了这些校验：几何构造用两套完全不同的方法互对（解析构造 vs 计算几何库
的精确求交），离散化用加密 $3$ 倍复算，文献闭式用数值特征分解对账，最优解用两条独立搜索
路线互认，对称性用镜像点核对，产物用逐字节比对。

\begin{table}[H]
  \centering
  \caption{自校验清单（每次运行自动执行）}
  \label{tab:checks}
  \begin{tabular}{p{0.52\textwidth}p{0.40\textwidth}}
    \toprule
    \textbf{校验项} & \textbf{结果} \\
    \midrule
    解析构造 vs 精确求交（$400$ 例，有效 $213$ 例） & 最大相对偏差 $1.2\times10^{-14}$ \\
    一阶式 vs 精确构造（$\gamma\in[30^\circ,120^\circ]$，$120$ 例） & 中位 $0.078\%$，$98\%$ 的样本 $<1\%$ \\
    离散化加密 $3$ 倍（$d$、$\delta_1$、$\delta_2$ 各 $3$ 倍）复算 $J^*$ & 相对偏差 $6.4\times10^{-15}$ \\
    最坏情形 vs 问题一精确区域口径 & 相对偏差 $0$ \\
    最坏情形是否被圆域截断 & 未截断，顶点数 $4$ \\
    可行域边界的最坏距离 & $\in[1000.0,1000.0]$ m（全部舍入误差内成立） \\
    候选弧带径向连续性 & 每个方位角均连续，断口 $0$ 处 \\
    镜像对称性（镜像点的代价） & 相对偏差 $\le1\times10^{-14}$ \\
    文献闭式 vs 数值特征分解（$53$ 例） & 最大相对偏差 $1.4\times10^{-12}$ \\
    全局认证（$5$ m 网格快筛 $+$ 精确复核 $+$ 细化） & $J=134.0831$ m，与粗搜解相隔 $8.49$ m（未漏盆地） \\
    确定性（两次运行的六个产物） & 逐字节一致 \\
    \bottomrule
  \end{tabular}
\end{table}

\subsection{模型优点}

本模型的第一个优点是\textbf{保证型}：不假设任何误差分布，"第二次一定听得到"与"最坏定位
直径"都是硬保证，直接对应 $20$ m 清除阈值这种硬判据，不会因为分布假设不成立而失效。第二个
优点是\textbf{可解释}：最优解的结构由两条几何事实完全解释——$1/\sin\gamma$ 主导（交会角越大
越好）与两个约束同时顶到边界（基线不超过 $1005$ m、源到第二点的距离不超过 $1000$ m）——
并有一阶解析式与文献的几何稀释结论互相印证，因此结论可以被现场人员理解与记忆。第三个优点
是\textbf{双重认证}：精确构造路线与"文献判据快筛 $+$ 精确复核"路线彼此独立，结果仅相差
$8.5$ m 且后者略优，消除了"粗网格漏掉更好盆地"这一类网格依赖风险。第四个优点是
\textbf{结论稳健}：minimax、文献 GDOP 与期望精度三种准则给出的最优点相距不足 $20$ m、
最坏直径相差小于 $0.1\%$，说明推荐区域不依赖判据的选择。最后，全部结果\textbf{完全可复现}：
程序无随机数，六个产物两次运行逐字节一致，论文中的每个数字都能由结论与校验文件反向核对。

\subsection{模型局限}

\begin{enumerate}
  \item \textbf{最坏界偏保守。} 有界误差模型必须承担"两个 $\pm1^\circ$ 误差同向叠加"的
        最坏情形，因此定位区域比其他口径大约三成。若工程上允许按统计模型评估风险，可用
        期望口径替代；本文已验证两种口径给出的最优点几乎相同（表~\ref{tab:three}），
        所以保守性对\textbf{选点}的影响很小，对\textbf{报出的精度指标}影响较大。
  \item \textbf{不确定集依赖题设上界。} $d\le1500$ m 来自接收半径上限；若该上界估计偏大，
        可行域会偏小、最坏直径偏大（安全方向）；若第一轮能给出距离先验，结果会显著更好
        （§\ref{sec:sens}）。
  \item \textbf{圆域截断按保守处理。} 代价函数用未截断四边形的直径作为上界。在本文的
        最优点处截断并不发生（已验证），但若检测点靠近圆域边界，算法会略微偏好靠内的解，
        属于安全方向的偏差。
  \item \textbf{未纳入执行代价。} 本问只回答"第二点选在哪里"，不考虑机器狗到达该点的里程、
        耗时与地形；两瓣镜像解在精度上完全等价，实际取舍应由路径成本决定。
  \item \textbf{未利用测向误差的空间相关性。} 题面指出同一位置、同一频道重复测量的误差相同
        （系统误差），这意味着"同一地点反复测量"没有信息增益。本文的选点天然满足"两次测
        向来自不同地点"，与该性质相容，但未进一步把它作为约束显式建模。
\end{enumerate}

\subsection{推广}

本模型的核心结构是"不确定集（集员）$+$ 保证型可测性约束 $+$ 最坏情况定位直径最小化 $+$
适合度场"，它可以推广到若干相邻情形。若已掌握部分距离先验，只需把不确定集的径向上界替换为
先验值，算法与图不需要任何改动（§\ref{sec:sens} 已给出该情形的完整结果）。若需要第三次及
以后的测向，把前两次的楔形交作为新的不确定集、重复同一套流程即可选出下一点，代价函数只需把
单项楔形交换成多项楔形交。对于 TDOA、RSSI 等其他观测体制，同样可以写成"不确定集 $+$ 区域
直径"的形式，文献判据层（Fisher 信息 $/$ CRLB）也都有对应形式，本文"判据对照 $+$ 快筛认证"
的框架可以直接复用。若源在观测期间移动，则把式~\eqref{eq:J} 的 $\max$ 换成沿时间轨迹的
$\max$（鲁棒轨迹优化），模型结构不变。

\subsection{结论}

对本题给定的一次测向结果，第二个检测点应当：

\begin{enumerate}
  \item 选在距第一检测点\textbf{尽可能远}（$1000$ m 量级，几何上限 $1005$ m）且相对示向度
        方向\textbf{偏向约 $37^\circ$} 的位置，最优解为
        $S_2^*=(801,605)$ m（$r^*=1004$ m、$\varphi^*=+37.1^\circ$，镜像解等价）；
  \item 该点的最坏定位直径为 $134.08$ m，相比只做一次测向的 $1800$ m 缩小 $13.4$ 倍，
        意味着在最坏情况下也能满足 $20$ m 清除阈值所需的定位收敛要求（配合后续补测）；
  \item 若需要现场灵活性，可在两条镜像弧带中选择：
        $r\in[947.6,1004.5]$ m、$\varphi\in\theta_1\pm[25^\circ,40^\circ]$，
        合计 $0.02069$ km$^2$，其中每一点的适合度都可用 $F=J^*/J$ 或图~\ref{fig:suit}
        的等值线直接读出；
  \item 该结论对判据选择（最坏 $/$ 文献 GDOP $/$ 期望）不敏感，对检测点位置与示向度具有
        平移旋转不变性，稳健可靠。
\end{enumerate}
